% !TeX root = ../main.tex
\chapter{Architectural Patterns}
\label{ch:patterns}

\section{How the Choice Was Made}

Architectural structure follows from the non-functional requirements more
closely than from the functional ones. A functional requirement can usually be
met by many structures; a non-functional one tends to rule most of them out.
Four of ours account for most of the structure.

\begin{longtable}{L{0.14\textwidth}L{0.36\textwidth}L{0.42\textwidth}}
	\caption{Non-functional requirements that drove the structure}\label{tab:drivers}\\
	\toprule
	\textbf{ID} & \textbf{What it asks for} & \textbf{Structural consequence} \\
	\midrule
	\endfirsthead
	\toprule
	\textbf{ID} & \textbf{What it asks for} & \textbf{Structural consequence} \\
	\midrule
	\endhead

	\req{NFR-8} & A generated diagram is readable without the source beside
	it. The governing criterion. & Discarding detail in the diagram is a decision in its own
	right. This is the Abstraction subsystem. \\
	\addlinespace
	\req{NFR-2} & A failure in the tool must not terminate GENESYS. & The
	analysis cannot share a process with the host. Two processes, not one. \\
	\addlinespace
	\req{NFR-4} & The same inputs always produce the same output, including
	ordering. & Each stage is a function of its input rather than of
	accumulated state, which is what makes repeated runs comparable. \\
	\addlinespace
	\req{NFR-6} & The projects that do not reference the GENESYS assemblies
	build without a GENESYS installation. & The vendor dependency is confined
	to one project, so the boundary between subsystems is also a boundary
	between licensed and our unlicensed code. \\
	\bottomrule
\end{longtable}

Maintainability points the same way. It is best served by fine-grained,
self-contained components, and a translation decomposes naturally into stages
that each do one thing. That is the shape of \pattern{Pipe and
Filter}, which is the primary pattern of this system.

\section{Pipe and Filter: The Primary Pattern}
\label{sec:pipe-and-filter}

\umlfigure{pipeline}{The translation pipeline. Each box is a filter; each
labelled arrow is a pipe carrying the named format}{fig:pipeline}

A translation is the processing of data in separate stages, each stage
performing one transformation and passing the result on. Figure~\ref{fig:pipeline}
shows the code-to-activity direction. The Code Extractor turns source files into
syntax trees and a control flow graph, which is \req{CA-1}. The Abstraction
filter discards constructs below the level a systems engineer reads, which is
where \req{NFR-8} is enforced. The Model Mapper turns what survives into
entities and relationships, which is \req{CA-2}.

A compiler is the standard example of this pattern, and this tool is a compiler
in the ordinary sense: it has a front end that parses, an intermediate
representation, and a back end that emits. That it reads a model rather than a language on one side
does not change the shape.

\subsection{Why It Fits}

\begin{longtable}{L{0.34\textwidth}L{0.58\textwidth}}
	\caption{Pipe and Filter against this system}\label{tab:pf-fit}\\
	\toprule
	\textbf{Property of the pattern} & \textbf{How it applies here} \\
	\midrule
	\endfirsthead
	\toprule
	\textbf{Property of the pattern} & \textbf{How it applies here} \\
	\midrule
	\endhead

	Inputs processed in separate stages to produce related outputs &
	Four translations, each a chain of stages over the same material. The
	outputs are related by construction, since they are produced from one
	abstract syntax. \\
	\addlinespace
	Supports transformation reuse & The review and write sequence is one
	filter shared by all four translations, which is what makes \req{CM-3} and
	\req{NFR-3} hold uniformly rather than being re-implemented per direction.
	The Abstraction filter is likewise shared. \\
	\addlinespace
	Seperate transformations allow for straightforward evolution & Adding Visual
	Basic means one more filter at the source end and no change to anything
	downstream. This is the concrete reason the tool supports two languages
	without two translators. \\
	\addlinespace
	Can be sequential or concurrent & The two halves already run in separate
	processes, which the pattern accommodates without alteration. That is what
	lets \req{NFR-2} be satisfied without distorting the design. \\
	\bottomrule
\end{longtable}

\subsection{The Cost We Accept}

Sommerville states the cost of this pipe and filter architecture directly: the format for data
transfer has to be agreed between communicating transformations, and each
transformation must parse its input and unparse its output to that form. That
cost appears here in two places.

The agreed form is \code{Pentateuch.Absyn}, described in
Section~\ref{sec:canonical}. Agreeing it was not free, and it is the reason a
third project exists where two would otherwise do. The unresolved half of the
cost is recorded as risk \textbf{R-2}: the mechanism by which the two processes
exchange that form has not been chosen. The abstract syntax fixes \emph{what} is
exchanged, so only the transport is open, but the pattern is what makes the
question unavoidable.

The second cost is less obvious. Because each filter must produce the agreed
form, a construct the abstract syntax cannot express cannot be carried. Rather
than let such a construct be silently dropped, \req{CA-8} and \req{AS-3} require
the tool to report it. The pattern did not create those requirements, but it is
what makes them necessary rather than optional.

\section{Repository: The Architecture We Integrate With}
\label{sec:repository}

\umlfigure{repository-pattern}{GENESYS organised as a repository. Views do not
communicate with each other}{fig:repository}

A reasonable question about any system with a database at its centre is why it is not
simply a \pattern{Repository} architecture. The question should be answered
here, because there is a repository and it is central to everything the tool
does.

GENESYS is itself a \pattern{Repository} architecture, in the strict sense of
the pattern. All model data lives in one store, and the components around it,
the activity diagram view, the state diagram view, the hierarchy view, do not
communicate with one another. Each regenerates from the repository on demand.
Figure~\ref{fig:repository} shows the arrangement with our Model Mapper drawn as
one more client of the same store.

This has a direct and useful consequence. Because diagrams are views generated
from the repository rather than stored pictures, the tool creates entities and
relationships and produces no layout, geometry or formatting. That is
\textbf{C-6}, and \req{CM-5} is the requirement that follows from it. A
property of a system we do not own removes a subsystem from ours. We do not need a diagramming component, and the output of a
translation is complete the moment the entities exist.

\subsection{Why Our Own Architecture Is Not a Repository}

The pattern requires that components interact only through the repository and
not directly. Ours do the opposite. The Code Extractor hands syntax trees to the
Abstraction filter, which hands abstract syntax to the Model Mapper. That is
direct interaction between stages, and it is how the design works.

The repository is where our pipeline ends, not how it is organised internally.
Drawing the system as a \pattern{Repository} architecture would also lose the
property that matters most for \req{NFR-4}: in a repository architecture, the
order in which components run is not fixed by the structure, whereas determinism
of output requires exactly that it be fixed.

\section{Client-Server: Our Position Relative to GENESYS}

The tool reaches the model through the GENESYS Client API, authenticating
against a repository that may be local or remote and depending on the Sentinel
licensing components to do so. That is a \pattern{Client-Server} relationship,
and we sit firmly on the client side of it. The relevant consequences are
constraints rather than choices: \textbf{C-4} requires a licensed local
installation wherever the tool runs, and \textbf{C-2} records that the server's
REST interface is not available to us, so all access goes through the client
assemblies.

Sommerville lists performance depending on the network as a disadvantage of this
pattern, and it is the reason the depth control of \req{CA-3} exists. The amount
of model traffic a translation generates is bounded by the user rather than by
the size of the repository.

One stretch requirement would change our position. \req{FS-1} would have the
tool run as a background service started and controlled by a GENESYS script, at
which point the tool becomes a server with GENESYS as its client. The design
keeps the application separable from any interface placed in front of it so that
this does not require restructuring. It is constrained by \textbf{C-3} rather
than by the architecture.

\section{Layered: A Local Principle, Not the Architecture}
\label{sec:layered}

\umlfigure{layering}{Layering holds inside a process and does not hold across
them}{fig:layering}

\pattern{Layered} architecture organises a system into tiers, each providing
services to the tier above, with dependencies pointing one way. Part of our
system is organised this way and part of it is not, and the distinction matters,
because claiming the whole system as layered would be wrong.

Inside each process the arrangement is genuinely layered, as the left panel of
Figure~\ref{fig:layering} shows. An access tier wraps the vendor APIs, a
translation tier turns what they return into our own structures, and a policy
tier decides what is significant enough to appear in the output. Each tier uses
only the one below it. This is the arrangement Sommerville's advantage applies
to: the access tier could be replaced wholesale, which is what makes the
interchange-file fallback of risk \textbf{R-1} viable without touching anything
above it.

Across processes the arrangement is not layered. \code{Pentateuch.Genesys} and
\code{Pentateuch.Translator} both depend on \code{Pentateuch.Absyn} and neither
depends on the other. Two components depending on a third is a shared kernel,
not a layer stack: neither of the two sits above the other, and there is no
service relationship between them to invert. Calling it a two-layer architecture
would describe the dependency graph without describing the design.

The honest summary is that layering is a local organising principle used inside
each process, and the system's architecture is the pipeline described in
Section~\ref{sec:pipe-and-filter}.

\section{Summary of the Five}

\begin{longtable}{L{0.20\textwidth}L{0.17\textwidth}L{0.55\textwidth}}
	\caption{Verdict on each architectural pattern considered}\label{tab:verdicts}\\
	\toprule
	\textbf{Pattern} & \textbf{Verdict} & \textbf{Reason} \\
	\midrule
	\endfirsthead
	\toprule
	\textbf{Pattern} & \textbf{Verdict} & \textbf{Reason} \\
	\midrule
	\endhead

	\pattern{Pipe and Filter} & Adopted, primary & A translation is staged
	data transformation. Gives transformation reuse across all four
	translations and makes a second language additive. \\
	\addlinespace
	\pattern{Repository} & Describes GENESYS, not us & Our stages interact
	directly, which the pattern forbids. Its presence on the GENESYS side is
	what makes \req{CM-5} achievable. \\
	\addlinespace
	\pattern{Client-Server} & Adopted at the boundary & The Client API
	relationship, with the tool as client. \req{FS-1} would reverse the roles
	for a future interface. \\
	\addlinespace
	\pattern{Layered} & Adopted locally & Holds inside each process as access,
	translation and policy tiers. Does not describe the relationship between
	processes. \\
	\addlinespace
	\pattern{MVC} & Not applicable & The tool has one view and no interactive
	editing of model data, so the problem this pattern solves does not arise.
	It would become relevant only if \req{FS-2}, the status window, is built. \\
	\bottomrule
\end{longtable}
